
\section*{Ventana N68: cilindro intervalar}
El criterio correcto no es el extremo inferior:
\[
\left\lfloor\frac{1000^K N}{3^L}\right\rfloor=D.
\]
El criterio correcto es intersección de cilindros:
\[
\left[\frac N{3^L},\frac{N+1}{3^L}\right)
\cap
\left[\frac D{1000^K},\frac{D+1}{1000^K}\right)
\neq\varnothing.
\]
Equivalente:
\[
N1000^K<(D+1)3^L,
\qquad
(N+1)1000^K>D3^L.
\]
Esto formaliza:
\[
\boxed{0_{\rm borde}\neq9_{\rm activo}.}
\]
El tail/carry puede empujar el prefijo al cilindro decimal; mirar sólo el extremo inferior pierde información de borde.
Con este criterio, \(t=12\) queda corregido y todos los bloques reales \(t=5,\ldots,20\) entran en sus cilindros.
\[
\boxed{\text{poda decimal debe ser intervalar, no puntual.}}
\]
